\documentclass[11pt]{article}

\usepackage{amsmath,amssymb}
\usepackage{xurl}
\usepackage[hidelinks]{hyperref}
\setlength{\emergencystretch}{2em}

% Shared commands for notes in docs/paper-gaps/.
%
% Two halves.  The link commands below are project identity: OWNER/REPO is the
% placeholder that scripts/bootstrap_project.py replaces with this project's
% GitHub slug and Pages site.  Everything after them is NOTATION, and notation
% belongs to the source paper: the definitions here are an example set, and a
% project replaces them with the macros of its own paper's style file.  The one
% rule that never changes: a command defined here must mean exactly what the
% source paper's macro means, or a note silently says something else.

\newcommand{\Lean}{\textsc{Lean}}
\newcommand{\leanid}[1]{\textsf{#1}}
\newcommand{\ghissue}[1]{\href{https://github.com/OWNER/REPO/issues/#1}{GitHub issue \##1}}
\newcommand{\ghpr}[1]{\href{https://github.com/OWNER/REPO/pull/#1}{PR \##1}}
% Local issue tracker (issues/NNNN-*.md in this repository); no remote URL.
\newcommand{\localissue}[1]{issue \##1 (\path{issues/})}
\newcommand{\doclink}[1]{\href{https://OWNER.github.io/REPO/docs/#1.html}{\path{#1}}}

% --- Notation: replace with the source paper's own macros --------------------
% \F, \N, \C, \R, \tr, \ot, \eps, \E, \bx, \bu, \bv, \by

\newcommand{\F}{\mathbb{F}}
\newcommand{\N}{\mathbb{N}}
\newcommand{\C}{\mathbb{C}}
\newcommand{\R}{\mathbb{R}}
\newcommand{\tr}{\mathrm{tr}}
\newcommand{\eps}{\epsilon}
\newcommand{\ot}{\otimes}
\newcommand{\E}{\mathop{\bf E\/}}

% bold random variables
\newcommand{\bu}{\boldsymbol{u}}
\newcommand{\bv}{\boldsymbol{v}}
\newcommand{\bx}{{\boldsymbol{x}}}
\newcommand{\by}{\boldsymbol{y}}

% T1 so literal < and > in \texttt placeholders typeset as themselves
\usepackage[T1]{fontenc}

% bra-ket
\usepackage{braket}

% --- Example structured spaces ---
\newcommand{\polyfunc}[3]{\mathcal{P}(#1, #2, #3)}
\newcommand{\polymeas}[3]{\mathrm{PolyMeas}(#1, #2, #3)}
\newcommand{\polysub}[3]{\mathrm{PolySub}(#1, #2, #3)}

% Approximation relations (paper uses \approx_\eps and \simeq_\zeta)
\newcommand{\approxeps}{\approx_{\varepsilon}}
\newcommand{\simgeq}{\simeq_{\zeta}}

\renewcommand{\ghissue}[1]{%
  \href{https://github.com/Dengnifer/MIPStarRE-A/issues/#1}{GitHub issue \##1}}

\title{Simultaneous Polynomial Measurements Are Not Obtained Coordinatewise}
\author{}
\date{2026-09-05}

\begin{document}
\maketitle

\section{At a glance}

\begin{itemize}
  \item \textbf{Difficulty.} Easy for simultaneity parameter $k = 1$, hard
  for general $k$.  The coordinatewise route planned for the formalization
  applies the single-polynomial soundness theorem to each coordinate and
  combines the resulting measurements by the sandwich of the source's
  Lemma~\texttt{lem:ld-sandwich}; an explicit pair of measurements refutes
  every bound this route could produce.  The source instead cites a
  reduction of the $k$-tuple test to the single-polynomial test in $m + k$
  variables at degree $d + 1$; that reduction is now carried out in the
  formalization, in the individual-degree formulation, where the degree stays
  $d$.
  \item \textbf{Estimated weight.} The case $k = 1$ is closed by relabeling
  the coordinate conclusions (about 150 lines).  The combining reduction is
  complete: a combined strategy for the directly indexed game in dimension
  $m + k$ at degree $d$, a transfer of the success probability with the
  constant $10$, and the extraction of $k$ polynomials from one polynomial by
  linearity, in about a dozen new definitions and 6000 lines of formal proof.
  \item \textbf{Mathlib/project split.} The counterexample is finite Fourier
  analysis over $\F_{2^s}$ and a Gram-matrix estimate.  The reduction itself
  is almost entirely project-local (game constructions and the consistency
  calculus), with the Schwartz--Zippel lemma and multivariate polynomial
  algebra taken from Mathlib.
  \item \textbf{Key mathematical inputs.} Additive characters of $\F_q^2$
  and isotropic vectors of a symmetric pairing in characteristic two; the
  eigenvalue bound for nearly orthogonal unit vectors; the reduction proving
  Theorem~4.43 from Theorem~4.40 in \cite{Natarajan2018NEEXP}.  Tracked in
  \ghissue{134}.
\end{itemize}

\section{Key theorem forms}

The source theorem (\texttt{lem:ld-soundness}) states that for every
projective strategy passing the low individual degree game with parameters
$(q,m,d,k)$ with probability at least $1 - \eps$ there are polynomial-tuple
measurements $G^{\mathrm A}$ and $G^{\mathrm B}$, with outcomes
$(g_1,\dots,g_k)$ of individual degree at most $d$, such that on average over
a uniform point $u \in \F_q^m$
\[
  A^{\mathsf{Point},u}_{a} \ot I \simeq_{\delta_{\mathrm{ld}}}
    I \ot G^{\mathrm B}_{[\mathrm{eval}_u(\cdot) = a]},
  \qquad
  G^{\mathrm A}_{[\mathrm{eval}_u(\cdot) = a]} \ot I \simeq_{\delta_{\mathrm{ld}}}
    I \ot B^{\mathsf{Point},u}_{a},
  \qquad
  G^{\mathrm A}_{g} \ot I \simeq_{\delta_{\mathrm{ld}}} I \ot G^{\mathrm B}_{g},
  \tag{source form}
\]
where $a = (a_1,\dots,a_k)$ and
$\mathrm{eval}_u(g_1,\dots,g_k) = (g_1(u),\dots,g_k(u))$.

The formalization planned to derive this from the case $k = 1$ as follows.
Reading the $i$-th coordinate of every answer yields $k$ strategies for the
game with $k = 1$, each passing with probability at least $1 - 3\eps$; the
single-polynomial theorem gives projective measurements $G^{w,i}$ with
\[
  A^{\mathsf{Point},u}_{[a_i]} \ot I \simeq_{\delta}
    I \ot G^{\mathrm B,i}_{[\mathrm{eval}_u(\cdot) = a_i]},
  \qquad
  G^{\mathrm A,i}_{[\mathrm{eval}_u(\cdot) = a_i]} \ot I \simeq_{\delta}
    I \ot B^{\mathsf{Point},u}_{[a_i]},
  \qquad
  G^{\mathrm A,i}_{g} \ot I \simeq_{\delta} I \ot G^{\mathrm B,i}_{g},
  \tag{coordinate form}
\]
where $A^{\mathsf{Point},u}_{[a_i]}$ is the $i$-th marginal of the joint
point measurement; and the planned combination was the sandwich
\[
  C^{w}_{g_1,\dots,g_k}
    = G^{w,k}_{g_k} \cdots G^{w,2}_{g_2}\, G^{w,1}_{g_1}\,
      G^{w,2}_{g_2} \cdots G^{w,k}_{g_k},
  \qquad
  A^{\mathsf{Point},u}_{a} \ot I \simeq_{C k (\delta + md/q)^{1/2}}
    I \ot C^{\mathrm B}_{[\mathrm{eval}_u(\cdot) = a]},
  \tag{planned form}
\]
with a universal constant $C$, and likewise on the other side and for the
global relation.  We show that no statement of the planned form holds: the
coordinate form, even with $\delta = 0$ and even with both point
measurements and all $2k$ polynomial measurements available, does not
determine any polynomial-tuple measurement consistent with the joint point
measurement.

The formalization proves the source form for every $k$, for the directly
indexed low-degree game of \cite{gap:qpbt_ld-dimension-divisibility} and by the
source reduction rather than by the refuted route.  From a projective strategy for the
directly indexed low-degree game with parameters $(q,m,d,k)$ succeeding with
probability at least $1 - \eps$ it produces polynomial-tuple measurements
satisfying the three relations of the source form with the error of the low
individual degree theorem at the combined parameters $(q, m+k, d, 1)$, at the
auxiliary parameter $2560000\,(m+k)^{3}d$ and pass bound $30\eps$, increased by
the recovery loss $(m+k)d/q$ on the two point/polynomial relations;%
\footnote{The declaration is
\leanid{exists\_directSimultaneousPolynomialMeasurements\_combinedError} in
\doclink{MIPStarRE/QPBT/Combining/DirectLowDegree/Transport/Combining/SimultaneousGeneral.lean}.}
and, absorbing that error, the same three relations with the error function
$\delta_{\mathrm{ld}}(\eps,q,m,d,k)$ of the source statement at the universal
constants $a = 10^{23}$ and $b = 1/80000$.%
\footnote{The declaration is \leanid{exists\_direct\_ld\_soundness} in
\doclink{MIPStarRE/QPBT/Combining/DirectLowDegree/Soundness.lean}.  The case
$k = 1$, the only one instantiated by the Pauli basis test analysis, remains
available separately as
\leanid{exists\_directSimultaneousPolynomialMeasurements\_of\_k\_eq\_one}
in \doclink{MIPStarRE/QPBT/Combining/DirectLowDegree/Transport/Simultaneous.lean};
the instantiation with $k = 1$ at dimension $2m+2$ is in the proof of
\texttt{lem:qld-4-7}, \path{blueprint/src/chapter/ch15_qpbt_combining.tex}.}

The route is the source's: the combined strategy answers each question of the
game at the combined parameters by measuring one question of the original
strategy and relabeling its outcome, its success probability is at least
$1 - 10\eps$, the case $k = 1$ is applied once to it, and the $k$ polynomials
are recovered from the single outcome polynomial by exact linearity in the
fresh coordinates.  Two quantitative points differ from the account of the
source assertion below, and both are forced by the individual-degree
formulation of the game.  The degree does not rise to $d + 1$: the combined
polynomial $(u,\alpha) \mapsto \sum_r \alpha_r g_r(u)$ has individual degree
at most $d$ in each of the $m + k$ variables, so the theorem is applied at
degree $d$ and its error function is evaluated there.  Correspondingly the
recovery cost is not $(d+1)/q$ taken twice but $(m+k)d/q$: the exact-linearity
test on the point space contributes $md/q$, and the collision of two distinct
tuples at the sampled point contributes $kd/q$.  A reader comparing the two
accounts should read the displayed recovery term of the source assertion as
$(m+k)d/q$ throughout.

\section{The source assertion}

The proof of \texttt{lem:ld-soundness} in
\cite[lines~413--458]{Ji2020MIPStar} treats $k = 1$ by an identification
with the tensor-code test of \cite{Ji2021Quantum} (the identification is
examined in \path{docs/paper-gaps/qpbt_ld-dimension-divisibility.tex}) and
then extends to general $k$ by reference to the derivation of Theorem~4.43
from Theorem~4.40 in \cite[lines~1409--1503]{Natarajan2018NEEXP}.  That
derivation is not coordinatewise.  It builds two combined provers who play
the single-polynomial test in $m + k$ variables at degree $d + 1$: on a
point $(\alpha, u)$ the combined prover asks the original prover for the $k$
values $(b_1,\dots,b_k)$ at $u$ and returns $\sum_i \alpha_i b_i$, and on a
line it returns the corresponding combination of the $k$ answer polynomials.
The degree rises by one because the combination
$(\alpha, u) \mapsto \sum_i \alpha_i g_i(u)$ is linear in $\alpha$ and of
degree $d$ in $u$, so its restriction to a queried subspace has degree
$d + 1$; the single-polynomial theorem is therefore applied at degree $d + 1$
and its error function is evaluated there.  That theorem then yields one
measurement with outcomes $g(\alpha,u)$, and the $k$ polynomials are
recovered from $g$ because $g$ is exactly linear in $\alpha$ with high
probability, by the Schwartz--Zippel lemma; applied at degree $d + 1$, that
lemma contributes the term $(d+1)/q$ twice in the recovery.  The
simultaneity of the $k$ outputs is thus inherited from the joint answers of
the original provers through a single application of the $k = 1$ theorem.

\section{The counterexample}

Let $q = 2^s$, let $d = 1$, and let $\Gamma$ be the set of univariate
polynomials $g = g_0 + g_1 x$ over $\F_q$, so $|\Gamma| = N = q^2$.  Two
distinct elements of $\Gamma$ agree at a uniformly random point of $\F_q$
with probability at most $1/q$.  Bob's Hilbert space is $\C^{\Gamma}$ with
standard basis $\{\ket{g}\}$; Alice's is a second copy, and the state is
the maximally entangled state
$\ket{\psi} = N^{-1/2} \sum_{g \in \Gamma} \ket{g}\ket{g}$, for which
$(X \ot I)\ket{\psi} = (I \ot X^{\mathsf T})\ket{\psi}$ and
$\bra{\psi} X \ot Y \ket{\psi} = N^{-1}\tr(X Y^{\mathsf T})$.

Equip $\Gamma \cong \F_q^2$ with the symmetric pairing
$\langle h, g\rangle = h_0 g_1 + h_1 g_0$ and the additive character
$\chi(t) = (-1)^{\mathrm{tr}(t)}$, where $\mathrm{tr}$ is the absolute
trace of $\F_q$; the pairing is nondegenerate, so
$\phi_h = N^{-1/2} \sum_{g} \chi(\langle h, g\rangle)\ket{g}$, $h \in \Gamma$,
is an orthonormal basis.  Evaluation at $u \in \F_q$ is represented by the
vector $v_u = u + x$, since $\langle h, v_u\rangle = h_0 + h_1 u = h(u)$, and
$\langle v_u, v_u\rangle = 2u = 0$ in characteristic two.  Bob's two
coordinate measurements are
\[
  G^{1}_h = \ket{\phi_h}\bra{\phi_h},
  \qquad
  G^{2}_g = \ket{g}\bra{g},
  \qquad
  Q^{1,u}_{[a]} = \sum_{h(u) = a} G^1_h,
  \qquad
  Q^{2,u}_{[b]} = \sum_{g(u) = b} G^2_g .
\]
Writing $H_u = \{h : h(u) = 0\} = \F_q v_u$ and summing the character over
the coset $\{h : h(u) = a\}$ gives the matrix elements
\[
  \bra{g'} Q^{1,u}_{[a]} \ket{g}
    = \frac1N \sum_{h(u) = a} \chi(\langle h, g - g'\rangle)
    = \frac{1}{q}\, \chi(\langle h_a, g - g'\rangle)\,
      \mathbf 1[\,g(u) = g'(u)\,],
\]
where $h_a$ is any polynomial with $h_a(u) = a$: the sum over
$H_u = \F_q v_u$ vanishes unless $\langle v_u, g - g'\rangle = 0$, that is,
unless $g(u) = g'(u)$, because $v_u$ is isotropic.  Hence $Q^{1,u}_{[a]}$
is block diagonal for the partition of $\Gamma$ by the value at $u$, it
commutes with every $Q^{2,u}_{[b]}$, its diagonal entries all equal $1/q$,
and each product $\Pi^u_{a,b} = Q^{1,u}_{[a]} Q^{2,u}_{[b]}$ is the rank-one
projector onto the unit vector
$\xi^u_{a,b} = q^{-1/2} \sum_{g(u) = b} \chi(\langle h_a, g\rangle)\ket{g}$.

Let both players answer the point question $u$ with the joint projective
measurement $\{\Pi^u_{a,b}\}_{a,b}$, which is real symmetric, and let Alice's
coordinate polynomial measurements be the transposes $(G^{i})^{\mathsf T}$.
On $\ket{\psi}$ every relation of the coordinate form holds with error
$0$: the marginals of the joint point measurement are $Q^{1,u}_{[a]}$ and
$Q^{2,u}_{[b]}$, which are exactly the evaluated coordinate measurements, and
$(G^i)^{\mathsf T} \ot I$ agrees with $I \ot G^i$ on $\ket{\psi}$.  The
collision probability is $1/q$.  The sandwich on Bob's side, evaluated at
$u$, is
\[
  C^{\mathrm B}_{[\mathrm{eval}_u(\cdot) = (a,b)]}
    = \sum_{g(u) = b} \ket{g}\bra{g}\, Q^{1,u}_{[a]}\, \ket{g}\bra{g}
    = \frac1q\, Q^{2,u}_{[b]},
\]
and its consistency with the joint point measurement is
\[
  \sum_{a,b} \bra{\psi} \Pi^u_{a,b} \ot C^{\mathrm B}_{[\mathrm{eval}_u(\cdot) = (a,b)]}
    \ket{\psi}
    = \frac{1}{qN} \sum_{a,b} \tr\bigl(Q^{1,u}_{[a]} Q^{2,u}_{[b]}\bigr)
    = \frac1q ,
\]
so the consistency defect is $1 - 1/q$ for every $u$, whereas the planned
form promises at most $2C\,(1/q)^{1/2}$.  No universal constant $C$ survives
$q \to \infty$; with the constant $C = 8$ established for the structured
sandwich, the planned bound already fails at $q = 2^{9}$.  The same
computation with the roles of the players exchanged refutes the planned
bound for Alice's sandwich against Bob's point measurement, and
$\sum_{g} \bra{\psi} (C^{\mathrm B}_g)^{\mathsf T} \ot C^{\mathrm B}_g \ket{\psi}
= N^{-1} \sum_{g_1,g_2} |\langle g_2 | \phi_{g_1}\rangle|^4 = 1/N$ refutes it
for the global relation.

The obstruction is not specific to the sandwich.  Let
$\{B_{g_1,g_2}\}$ be any POVM on Bob's space with polynomial-pair outcomes,
and write $B_{g_1,g_2} = N p_{g_1,g_2}\,\rho_{g_1,g_2}$ with
$p$ a probability distribution and $\rho$ density matrices.  Its
consistency with the joint point measurement is
\[
  \E_{u} \sum_{g_1,g_2} p_{g_1,g_2}\,
    \bra{\xi^u_{g_1(u),g_2(u)}}\, \rho_{g_1,g_2}\, \ket{\xi^u_{g_1(u),g_2(u)}}
  \le \max_{g_1,g_2}\,
    \lambda_{\max}\Bigl(\E_{u}\, \ket{\xi^u_{g_1(u),g_2(u)}}\!\bra{\xi^u_{g_1(u),g_2(u)}}\Bigr).
\]
For $u \ne u'$ the supports $\{g : g(u) = g_2(u)\}$ and
$\{g : g(u') = g_2(u')\}$ meet exactly in $\{g_2\}$, so the $q$ unit vectors
$\xi^u_{g_1(u),g_2(u)}$ have pairwise inner products of modulus $1/q$, and
the largest eigenvalue of their Gram matrix is at most $1 + (q-1)/q < 2$.
Hence the consistency of every polynomial-pair POVM with the joint point
measurement is below $2/q$, on either side.  Consequently no argument that
uses only the coordinate form --- with any error, any number of applications
of the pasting or sandwich facts, and either player as reference --- can
produce the source form for $k \ge 2$.

Two remarks.  First, the pair $(\ket{\psi}, \Pi, G^1, G^2)$ is not a
strategy for the low-degree game passing with high probability: by the same
Gram-matrix bound no line measurement with polynomial-pair outcomes is
consistent with the point measurements, so the source theorem is not
contradicted; what fails is the coordinatewise route.  Second, the pasting
fact (Fact~4.35 of \cite{Natarajan2018NEEXP}, imported as
\texttt{lem:pasting}) does not repair the route: it requires the two
coordinates to be evaluated at points $y_1, y_2$ whose conditional
distribution supports a collision bound $\eta < 1$, whereas the low-degree
game evaluates all coordinates at the same point, for which the conditional
collision probability of any two polynomials agreeing there is $1$.

\section{Comparison with the blueprint and formalization}

The blueprint proof of \texttt{lem:ld-soundness} records the source route,
including the extension from $k = 1$ to general $k$ by the derivation of
Theorem~4.43 in \cite{Natarajan2018NEEXP}, and lists the game
correspondence and the parameter bound as open obligations.  The present
note added a third: the extension is the combining reduction and must be
formalized as such.  That third obligation is now met.  The blueprint records
the reduction in \texttt{prop:ld-simultaneous-general-k} and the nodes
preceding it, and the directly indexed game and its transport prove the source
form from the low individual degree theorem for every $k$.  The game
correspondence and the parameter bound of the import remain as they were.  The
combination of the coordinate measurements by the sandwich is defined and is a
POVM, but it carries no consistency estimate, and by the computation above it
cannot; nothing depends on it.

\section{Conclusion}

The source statement is not contradicted.  Its proof for general $k$ is a
reduction to a single application of the $k = 1$ theorem in $m + k$
variables, and the formalization follows that reduction; a
coordinatewise application of the $k = 1$ theorem followed by any
combination of the coordinate measurements provably cannot yield the
theorem.  The application in the analysis of the Pauli basis test uses
$k = 1$ only, for which the source form is proved from the low individual
degree theorem directly; for general $k$ it is proved by the reduction, for
the directly indexed game.

\bibliographystyle{alpha}
\bibliography{references}

\end{document}
